% Derived from research/memory/exact_query_memory.tex; research original preserved.
\section{Exact empirical residual-query memory}
\label{sec:memory}
For the remainder of this section let
\[
 p=\binom{d+1}{2},\qquad q=\binom{d+3}{4},\qquad
 D_0=p+q,\qquad D=2p+q.
\]
These count symmetric second moments, symmetric fourth moments,
and the label--quadratic cross moment. The continuous-dimension
result concerns a fixed known sample count $n\ge D_0$.
\subsection{Definition, query equivalence and a matching summary}

Fix $d,n\ge1$ and a known sample count $n$.  Data are an ordered tuple
$z=((x_i,y_i))_{i=1}^n\in(\R^d\times\R)^n$.  Permutation invariance is
not required of a competing summary; allowing order only strengthens the
lower bound. Set
\begin{align*}
 H(x)&=xx^\top-I_d,\qquad q_B(x)=\tr(BH(x)),\\
 G_z(B,c)&=\frac1{2n}\sum_{i=1}^n
       (y_i-q_B(x_i)-c)H(x_i),\quad B\in\Sym_d,\ c\in\R.
\end{align*}
A summary $S_n:z\mapsto\R^m$ is \emph{universally query-sufficient} if
there is a decoder $\Gamma_n$ such that
\[
 \forall z\ \forall B\in\Sym_d\ \forall c\in\R:
 \quad \Gamma_n(S_n(z),B,c)=G_z(B,c).
 \tag{U}
\]
No regularity of $\Gamma_n$ is needed for the continuous lower bound.
The decoder has no additional data-dependent side information unless
explicitly stated; all retained model/optimizer/replay information that
it accesses must be counted in the summary. A known sample count is free
in this fixed-$n$ model. A streaming implementation stores the count too.

Write $\langle a\rangle_z=n^{-1}\sum_i a(x_i,y_i)$ and define
\[
 S_y=\tfrac12\langle yH\rangle_z,\qquad
 M_{2,ij}=\langle x_ix_j\rangle_z,\qquad
 M_{4,ijkl}=\langle x_ix_jx_kx_l\rangle_z.
\]
The subscripts on $M_2,M_4$ are coordinate indices rather than sample
indices. Store only distinct symmetric coordinates: $i\le j$ for $M_2$
and $i\le j\le k\le l$ for $M_4$. Let
\[
 F_n(z)=M(z)=(S_y,M_2,M_4)\in\R^D.
\]
The order-four tensor has $q$ distinct entries, not $p(p+1)/2$ independent
entries. Products of centered quadratic features have polynomial
relations, and known total mass is not an additional unknown coordinate.

\begin{theorem}[Canonical query information]
\label{thm:memory-canonical}
For any two datasets $z,z'$ of the same known size, the following are
equivalent:
\begin{enumerate}
 \item $G_z(B,c)=G_{z'}(B,c)$ for every symmetric $B$ and every $c$;
 \item $(S_y,\overline H,T)_z=(S_y,\overline H,T)_{z'}$, where
 $\overline H=\langle H\rangle$ and
 $T(B)=\tfrac12\langle\langle B,H\rangle H\rangle$;
 \item $M(z)=M(z')$.
\end{enumerate}
Consequently $M$ answers every query, and every sufficient summary
permits recovery of $M$. This is minimality for a deterministic
query family, not a statement of statistical sufficiency for an unknown
parameter.
\end{theorem}
\begin{proof}
Direct expansion gives
\[
 G_z(B,c)=S_y-T(B)-\tfrac12c\overline H. \tag{M1}
\]
Thus the value at $(0,0)$ determines $S_y$, its difference from $(0,1)$
determines $\overline H$, and its differences from queries at a basis of
$\Sym_d$ determine $T$. Conversely these coefficients determine every
query.

For $k\le l$, put $F^{kk}=e_ke_k^\top$ and
$F^{kl}=(e_ke_l^\top+e_le_k^\top)/2$ if $k<l$, so that
$\langle F^{kl},H\rangle=H_{kl}$. Define
\[
 K_{ij,kl}=\langle H_{ij}H_{kl}\rangle
             =2[T(F^{kl})]_{ij}.
\]
The reconstruction formulas are
\begin{align}
 M_{2,ij}&=\overline H_{ij}+\delta_{ij},\tag{M2}\\
 M_{4,ijkl}&=K_{ij,kl}
       +\delta_{ij}M_{2,kl}+\delta_{kl}M_{2,ij}
       -\delta_{ij}\delta_{kl}.\tag{M3}
\end{align}
Conversely (M3) reconstructs $K$ from $M_2,M_4$, and hence reconstructs
$T$. All index permutations are determined by tensor symmetry.
If two histories have the same sufficient summary, (U) makes all their
queries agree, so the reconstructed $M$ agrees too. Therefore a map
$h$ on the summary image satisfies $M=h\circ S_n$.
\end{proof}

The $p+2$ fixed queries
\[
 \mathcal Q=\{(0,0),(0,1)\}\cup
              \{(F^{kl},0):1\le k\le l\le d\}\tag{M4}
\]
already recover all coefficients. This reduces any exact universal
adaptive-query claim to a finite nonadaptive query family. Once $M$ is
stored, correctness is pathwise and holds for arbitrarily data-dependent
choices of $B,c$; independence or generalization assumptions are unnecessary.
If only one prespecified algorithm trajectory must be supported, (U)
is a stronger requirement and the lower bound need not apply.

The summary is additive: store sums of $yH/2$, all degree-two input
monomials, and all degree-four input monomials, then divide by $n$.
This uses $D$ ideal real coordinates and a count, independently of stream
length. It is an algebraic exact-arithmetic representation; ordinary
floating-point accumulation has rounding error. Nothing here identifies
the real-coordinate count with a physical byte count.

\subsection{Continuous summaries on fixed unweighted samples}

Let $f(x)\in\R^{D_0}$ list all degree-two and degree-four monomials.
The functions $1,f_1,\ldots,f_{D_0}$ are linearly independent. The
additional $p$ functions $yH_{ij}(x)/2$ are independent modulo this
input-only span. Indeed, the coefficient of $y$ in a linear relation
would say $x^\top Ax-\tr(A)\equiv0$, forcing $A=0$; independence of
different input monomials then finishes the argument. Thus the affine
span of the single-example canonical feature vector is all of $\R^D$.

\begin{lemma}[A full-dimensional unweighted moment image]
\label{lem:memory-fullrank}
For every fixed $n\ge D_0$, there is a dataset $z_0$ at which the
polynomial moment map $F_n:(\R^{d+1})^n\to\R^D$ has differential of
rank $D$. Consequently its image contains a nonempty open subset of
$\R^D$, and there is a continuous local section $s:U\to(\R^{d+1})^n$
on some nonempty open $U\subset\R^D$ such that $F_n(s(u))=u$.
\end{lemma}
\begin{proof}
First consider $A_n(x_1,\ldots,x_n)=n^{-1}\sum_i f(x_i)$.
The vectors $\partial_a f(x)$, over all $x\in\R^d$ and coordinates
$a$, span $\R^{D_0}$. Otherwise a nonzero $v$ annihilates them all,
so every partial derivative of the polynomial $v^\top f$ vanishes.
That polynomial is constant on $\R^d$, contradicting independence
of $1,f_1,\ldots,f_{D_0}$.

Select $D_0$ such derivative vectors forming a basis, and use their
locations as $D_0$ sample points. Their selected derivative columns in
$DA_{D_0}$ are independent. Pad the configuration when $n>D_0$.
Hence some $D_0$-minor of $DA_n$ is a nonzero polynomial in the sample
coordinates. The configurations where it is nonzero form an open dense
set $\mathcal U$.

Independence of the $p$ polynomials $H_{ij}$ implies that there are
$p$ points for which the $p\times p$ evaluation matrix has rank $p$:
if their value vectors never spanned $\R^p$, a nontrivial linear
combination of the functions would vanish everywhere. Since $n\ge
D_0\ge p$, the $p\times n$ matrix with columns $H(x_i)$ in symmetric
coordinates has rank $p$ on another open dense set $\mathcal V$.
Both density assertions follow because a nonzero real polynomial
cannot vanish on an open set. The finite intersection
$\mathcal U\cap\mathcal V$ is nonempty.

Choose an input configuration there and set every label to zero.
At this dataset, the derivative of $S_y$ with respect to input
coordinates is zero, while its label-derivative columns are
$H(x_i)/(2n)$, of rank $p$. The derivative of $(M_2,M_4)$ with respect
to inputs has rank $D_0$ and its label derivatives vanish. Thus the
full differential has rank $p+D_0=D$. Choose $D$ independent scalar
input/label coordinate directions and freeze the others. The inverse
function theorem applied to the resulting square map supplies $U$
and its smooth, hence continuous, local section.
\end{proof}

\begin{theorem}[Tight continuous dimension for sufficiently many samples]
\label{thm:memory-dimension}
Fix $n\ge D_0$. Any continuous universally query-sufficient summary
$S_n:(\R^{d+1})^n\to\R^m$ must have $m\ge D$. This bound is attained
by $M$. If $S_y$ is supplied separately and an additional continuous
summary $A_n$ has values in $\R^m$, exact universal decoding from
$(S_y,A_n)$ requires and permits $m\ge D_0$.
\end{theorem}
\begin{proof}
On the local section in the lemma, $S_n\circ s:U\to\R^m$ is
continuous. It is injective: equal summary values imply equal query
coefficients by the first theorem, and $F_n(s(u))=u$.
If $m<D$, compose this injection with the usual coordinate embedding
$\R^m\hookrightarrow\R^D$. Brouwer's invariance of domain says its
image is open in $\R^D$, but it lies in an $m$-dimensional coordinate
subspace with empty interior. This is a contradiction.

For separate $S_y$, the combined map $(S_y,A_n)$ is a continuous
summary of dimension $p+m$, so $p+m\ge D$. Storing $M_2,M_4$
attains $m=D_0$. Neither proof assumes continuity of the decoder.
\end{proof}

\paragraph{Finite sample qualification.}
The sufficient threshold $n\ge D_0$ is not claimed to be the smallest
threshold. For general $n$, let
\[
 r_n=\max_z\rank DF_n(z).
\]
The same independent-minor argument and a local parameter slice imply
that every continuous universal summary has dimension at least $r_n$.
Moreover $r_n\le\min\{D,n(d+1)\}$ and storing the ordered raw data
uses $n(d+1)$ continuous coordinates. Thus a $D$ lower bound cannot
hold when $n(d+1)<D$. At $n=1$, for example, storing $(x,y)$ uses
only $d+1$ coordinates. These bounds do not determine the exact
minimum for every small $n$; a sharper small-$n$ classification can
involve the geometry of finite-atomic polynomial moment varieties.

\paragraph{Quantifiers and model restrictions.}
The theorem assumes correctness at every dataset in the stated
unrestricted domain. It is not automatically a lower bound for data
on a fixed noiseless teacher manifold, for a discrete alphabet, or
for almost-sure correctness with an arbitrary discontinuous decoder.
Even topological full support of a sampling distribution does not turn
an almost-sure factorization through a discontinuous decoder into a
pointwise factorization. If both encoder and the finitely many
decoders in (M4) are continuous, equality almost surely under a
full-support law extends by continuity to every dataset, so the
theorem then applies. A fixed shared random seed also reduces to this
deterministic theorem only if that seed's map has universal correctness.

\subsection{Linear measurements and why additivity alone is insufficient}

The domain in this section is all finitely supported probability
measures $\mu$ on $\R^{d+1}$, with arbitrary nonnegative real weights,
not merely equal weights. Let $\phi_1,\ldots,\phi_m$ be any real-valued
functions and store their linear measurements $L(\mu)=(\int\phi_jd\mu)_j$.
Here ``linear'' means linear in the measure, not linear in $x$ or $y$.

\begin{proposition}[Exact linear-measurement criterion]
\label{prop:memory-linear}
For a target function $f$, the scalar $\int f\,d\mu$ is determined
by $L(\mu)$ for every finitely supported probability measure if and
only if
\[
 f\in\operatorname{span}_{\R}\{1,\phi_1,\ldots,\phi_m\}.
 \tag{M5}
\]
Therefore universal exact quadratic residual queries require
$m\ge D$ in this measurement model, regardless of decoder regularity.
The canonical measurements attain equality.
\end{proposition}
\begin{proof}
The ``if'' implication follows by integrating a linear representation,
using known total mass one. For the converse, put
$v(z)=(1,\phi_1(z),\ldots,\phi_m(z))$. If every finite real relation
$\sum_i a_iv(z_i)=0$ also implied $\sum_i a_if(z_i)=0$, then
$\ell(\sum_i a_iv(z_i))=\sum_i a_if(z_i)$ would be a well-defined
linear functional on the span of the $v(z)$. Extend it to $\R^{m+1}$;
then $f(z)=\ell(v(z))$, establishing (M5).
If (M5) fails, there is instead a finite relation with
$\sum_i a_i=0$ and $\sum_i a_if(z_i)\ne0$. Split the coefficients
into positive and negative parts. Their common positive mass $a$
normalizes them into two finitely supported probability measures
$\mu_+,\mu_-$ with equal $L$ but unequal integrals of $f$.
No decoder can distinguish them.

Apply this criterion to the $D$ canonical feature functions. Their
classes modulo constants are linearly independent, as proved above,
so the span of $m$ measurement classes must have dimension at least
$D$.
\end{proof}

An equivalent finite-support kernel statement is
$\ker L\cap\ker(\text{mass})\subseteq\ker M$: every signed measure
annihilated by the stored measurements and mass must annihilate all
query coefficients. This is ordinary linear algebra. It is not a
separate streaming or information-theoretic principle.

\paragraph{Two models that must not be conflated.}
For equal-weight samples of a fixed size, arbitrary discontinuous
additive features do not satisfy the lower bound in the proposition.
Choose a Hamel basis of $\R$ over $\mathbb Q$ and assign each possible
record $z\in\R^{d+1}$ to a distinct basis element $b_z$; both sets
have cardinality continuum. The scalar sum $\sum_i b_{z_i}$ uniquely
specifies every multiplicity in the finite multiset, by rational
linear independence. Thus one additive exact real, together with the
known count, stores all empirical query information. This encoding is
not a continuous finite-precision algorithm. Its existence shows why
one cannot infer a lower bound from ``additive'' plus a nonlinear
decoder alone. Arbitrary-real-weight mixtures defeat this example
because real, rather than rational, linear relations are then allowed.
Likewise, without any encoder regularity, the cardinality equality
$|\R^D|=|\R|$ permits storing the whole moment vector in one ideal
real. Even measurability alone is not a replacement for continuity:
standard digit encodings can be made Borel measurable.
This discontinuous-sum loophole is established prior theory: see
Wagstaff et al., Theorem 3.2 and Supplement B.2\cite{wagstaff}.

If a decoder is required to be affine in the stored feature averages,
the span criterion and $D$ lower bound follow already from testing
single records (or a dataset consisting of $n$ copies of one record).
This is a different, stronger restriction on decoding.

\subsection{A finite-bit lower bound with explicit error metric}

Let $\|A\|_{\max}=\max_{ij}|A_{ij}|$. Define simultaneous query error
on the fixed finite family (M4) by
\[
 e_z=\max_{(B,c)\in\mathcal Q}
       \|\widehat G_z(B,c)-G_z(B,c)\|_{\max}.\tag{M6}
\]
This metric is specified because an unqualified phrase such as
``$\varepsilon$-accurate residual queries'' does not define a memory
problem. All matrices $F^{kl}$ have Frobenius norm at most one, so
the query family uses bounded coefficients.

\begin{lemma}[Stable recovery from the finite query family]
\label{lem:memory-stability}
If $e_z\le\varepsilon$, the formulas in the first theorem reconstruct
a vector $\widehat M$ satisfying
$\|\widehat M-M(z)\|_\infty\le12\varepsilon$.
\end{lemma}
\begin{proof}
The errors in $S_y=G(0,0)$, $\overline H=2(G(0,0)-G(0,1))$,
and $K_{ij,kl}=2(G(0,0)-G(F^{kl},0))_{ij}$ are at most
$\varepsilon$, $4\varepsilon$, and $4\varepsilon$, respectively.
Equation (M2) has error at most $4\varepsilon$. Equation (M3) adds at
most two such errors to the error of $K$, giving $12\varepsilon$.
To select one value for each fourth-order coordinate, fix any
deterministic pair grouping in (M3); approximate reconstructions need
not otherwise respect redundant tensor identities.
\end{proof}

\begin{theorem}[Local worst-case bit bound]
\label{thm:memory-bits}
Fix $d\ge1$ and $n\ge D_0$. There exist a moment vector $m_0$,
$r>0$, and a continuous section $s$ on a neighborhood of
$C=m_0+[-r,r]^D$, with $F_n(s(m))=m$. Fix the compact family of
size-$n$ unweighted datasets
\[
 \mathcal K=s(C).
\]
Any summary with at most $2^b$ states and any decoder satisfying
(M6) with $e_z\le\varepsilon$ for every $z\in\mathcal K$ must obey,
for $\varepsilon>0$,
\[
 b\ \ge\ D\log_2\left(1+\left\lfloor
                   \frac{2r}{25\varepsilon}\right\rfloor\right).
 \tag{M7}
\]
If $S_y$ is supplied exactly as uncharged side information, the
number of additional bits has the same lower bound with $D_0$
replacing $D$. No finite number of states can be exact on all of
$\mathcal K$. These claims require correctness only on this compact
family, not on all real-valued datasets.
\end{theorem}
\begin{proof}
The local section gives an open realizable moment set. Choose a
closed cube $C$ inside it; its continuous image $\mathcal K=s(C)$
is compact. In each coordinate place
$L=1+\lfloor2r/(25\varepsilon)\rfloor$ equally spaced grid values
with spacing $25\varepsilon$; choose one realizing dataset via the
section for each of its $N=L^D$ grid points. If two grid datasets
had the same stored state, the decoder would return the same fixed
query outputs for both. Recovering $M$ from these outputs would
approximate both moment vectors within $12\varepsilon$, so their
distance would be at most $24\varepsilon$, contradicting the grid
separation. Hence there are at least $N$ states, proving (M7).

For side information, fix the first $p$ coordinates of the cube
(the label moment) at their center, and vary only the remaining
$D_0$ coordinates. All datasets now have identical side information,
and the same proof applies. For zero error, even a line segment
inside either cube has infinitely many distinct exact query vectors,
so finite-state universal exact recovery is impossible.
\end{proof}

The compact-family restriction makes a fixed positive absolute error
meaningful. If one instead required the same guarantee over all
unbounded real-valued datasets, already $G_z(0,0)=S_y$ has unbounded
range, so no finite-state code could have a finite uniform absolute
error. Such a global impossibility would not be a useful quantitative
precision bound.

This is a packing/pigeonhole theorem, not an invocation of Fano's
inequality without a probability model. It has an optional precise
randomized version. Draw the dataset uniformly from the constructed
$N$-point packing. A randomized $b$-bit encoding/decoding scheme with
shared randomness independent of the dataset and simultaneous
success probability at least $1-\delta$ on each dataset satisfies
\[
 b\ge\log_2N+\log_2(1-\delta).\tag{M8}
\]
Condition on all random coins: each state can be successful for at
most one packed dataset, so at most $2^b$ of the $N$ datasets
succeed; average this inequality over the coins. This also covers
private coins by revealing all of them as free shared randomness.
A per-query success guarantee is different from simultaneous
success for (M4), and requires its own union-bound adjustment.

For fixed $d,n$ and the resulting local scale $r$, (M7) gives
$D\log_2(1/\varepsilon)-O_{d,n,r}(1)$ as $\varepsilon\downarrow0$.
The proof does not give a dimension-uniform lower bound on $r$, a
Gaussian-prior rate-distortion function, or a practical fixed-precision
word count. The same compact family has a matching-order elementary
upper bound. For $\eta=\varepsilon/3$, use a grid on each canonical
coordinate interval of length $2r$ with
$J=1+\lceil2r/\eta\rceil$ values and spacing at most $\eta$.
Encoding the nearest grid vector has coordinate error at most $\eta$
and uses no more than
\[
 b_{\rm upper}=\left\lceil D\log_2\left(
                   1+\left\lceil\frac{6r}{\varepsilon}\right\rceil
                              \right)\right\rceil
 \tag{M9}
\]
bits. In (M1)--(M3), the errors in $S_y$, $\overline H$, and $K$ are
at most $\eta$, $\eta$, and $3\eta$; hence every query in (M4) has
error at most $3\eta\le\varepsilon$. The grid vector need not
itself be a realizable empirical moment vector for this decoder to
work. If $S_y$ is supplied exactly, quantize only the $D_0$ input
moment coordinates and replace $D$ by $D_0$ in (M9). Thus on the
specified compact family, the lower and upper bounds have matching
$D\log_2(1/\varepsilon)$ leading order for fixed $d,n,r$.

In particular, neither the continuous exact theorem nor this local
high-accuracy packing result proves that useful approximate residual
queries require $\Omega(d^4)$ storage at a fixed practical tolerance.
Classical covariance sketches such as Frequent Directions explicitly
trade storage against covariance approximation error\cite{p7ghashami2015}; their
applicability and bounds for the lifted quadratic features require a
separate calculation. The local radius $r$ in (M7) must not be silently
treated as dimension independent.

\subsection{Concrete collisions and the Phase6 boundary}

\paragraph{The missing fourth moment is operational.}
For $d=1,n=2$ and both labels zero, take dataset A with $x=(-1,1)$
and dataset B with $x=(0,\sqrt2)$. Both have $S_y=0$ and
$M_2=1$ (equivalently $\overline H=0$). Dataset A has
$\langle H^2\rangle=0$, whereas dataset B has
$\langle H^2\rangle=1$. At $B=1,c=0$, their exact query values
are $0$ and $-1/2$. Thus even $S_y$ together with the empirical
second moment cannot answer all residual queries. Replicating each
record equally gives the same collision at any even sample size.

\paragraph{Implementation inspected.}
The Phase~6 source \texttt{code/frontier.py}\cite{phase6} updates
\[
 \texttt{moment}\mathrel{+}=
 \tfrac12\sum_{\text{new }i}y_i(x_ix_i^\top-I),\qquad
 \texttt{matrix}=\texttt{moment}/n-B.
\]
Thus the proposal uses $S_y-B$. The exact historical empirical
quantity is (M1). For Gaussian inputs, $\E T(B)=B$ for fixed
sample-independent $B$ and $\E\overline H=0$, but these are
population identities. They do not turn $S_y-B$ into the realized
empirical query. The Phase~5 and Phase~6 theory files\cite{phase5,phase6} already state
this distinction correctly; this note supplies qualified minimality
results that those files deliberately did not claim.

With the Phase6 value $d=12$,
\[
 p=78,\quad q=1365,\quad D_0=1443,\quad D=1521.
\]
The sufficient sample-size threshold in the continuous lower bound
is therefore $n\ge1443$. With 1024 observations per arrival it
applies from the second arrival if the query target is the whole
historical empirical dataset. This observation is conditional on
(U) and on what information the decoder can access. It is not a
claim that the implemented approximate controller violates a memory
constraint. A dense symmetric $p\times p$ Gram matrix plus
$\overline H$ would use $p(p+1)/2+p=3159$ stored scalar slots,
whereas the input polynomial moments use only $1443$; adding $S_y$
gives $3237$ versus $1521$. These compare representations, not bit
lower bounds or runtime. Known constants and exact identities explain
the redundancy.

\paragraph{Changing the requested output changes the target dimension.}
The matrix query $G(B,c)$ does not contain $\langle y\rangle$ or
$\langle y^2\rangle$. To reproduce the intercept gradient as well,
store $\langle y\rangle$, one more independent coordinate:
\[
 \langle y-q_B-c\rangle
 =\langle y\rangle-\langle B,\overline H\rangle-c.
\]
To reproduce the entire squared empirical loss, add
$\langle y^2\rangle$ too. The corresponding unrestricted moment
dimensions are $D+1$ and $D+2$. Exact matrix gradient queries
therefore do not alone reproduce factor Adam with its chosen
minibatches, replay selection, coordinatewise states, and clipping.

\paragraph{Removing the intercept is dimension dependent.}
One must not simply subtract $p$ coordinates when restricting queries
to $c=0$. For $d\ge3$, the centered Gram moments already determine
$\overline H$. The pointwise rank-one identities of $H+I$ imply
\begin{align*}
 \langle H_{ii}+H_{jj}\rangle
    &=\langle H_{ij}^2-H_{ii}H_{jj}\rangle-1
            &&(i\ne j),\\
 \langle H_{jk}\rangle
    &=\langle H_{ij}H_{ik}-H_{ii}H_{jk}\rangle
            &&(i,j,k\text{ distinct}).
\end{align*}
Three pairwise diagonal sums recover each diagonal entry, and the
second identity recovers each off-diagonal entry. Thus for $d\ge3$
the dimension remains $D$ even without $c$. For $d=1$, the only
input query coefficient is $\langle H^2\rangle$, so the total affine
moment dimension is two including $S_y$. For $d=2$, the six distinct
Gram polynomials are independent modulo constants: their degree-four
leading terms force all coefficients to vanish except possibly a
multiple of $H_{11}H_{22}-H_{12}^2$, and its degree-two part is
nonzero. Hence the total is $6+p=9$, compared with $D=11$ when all
intercepts are queried. These are affine moment-space counts; the
finite-$n$ qualifications still apply.

\subsection{Prior theory and defensible contribution}

\begin{itemize}
\item \textbf{Classical least squares.} In the feature vector
$h(x)=\operatorname{vec}_{\Sym}(xx^\top-I)$, equation (M1) is the
normal-equation residual $X^\top y-X^\top X\beta$, with the intercept
cross term. MIT's official notes give the ordinary normal equations
and their gradient derivation\cite{mit}. The matrix query identity,
Gram storage, and regression interpretation are established theory.
\item \textbf{Statistical sufficiency is a different quantifier.}
The official Berkeley notes define sufficiency through the
conditional distribution of data given a statistic, relative to a
specified statistical model, and minimal sufficiency through recovery
from any sufficient statistic\cite{berkeley}. This note instead
fixes a family of deterministic empirical queries. For a fixed design
Gaussian regression model, the design is conditioned upon and its
Gram matrix need not be charged as stored unknown information. This
is another reason not to borrow statistical minimal-sufficiency
claims as memory bounds here.
\item \textbf{Polynomial moments and kernel means.}
Sriperumbudur et al. study mean embeddings of measures and conditions
for full distributional injectivity\cite{kme}. Taking the finite
feature map to be our canonical polynomials makes $M$ its empirical
kernel mean. It distinguishes this query family but does not
determine the full data distribution. The finite-dimensional feature
interpretation, and the distinction between query equivalence and
a characteristic kernel, are established theory.
\item \textbf{Moment compression is classical.}
Bayer and Teichmann prove a version of Tchakaloff's theorem: finitely
many weighted evaluation nodes can reproduce polynomial integrals
through a fixed degree\cite{tchakaloff}. This concerns existence of
positive cubature representations, not a lower bound on physical
memory. It reinforces that finite polynomial moment representations
are old, and does not remove the coordinate or bit cost of describing
nodes and weights.
\item \textbf{Topology is classical.}
The only topological input is Brouwer's invariance of domain;
Tao's author-written exposition gives the theorem, its no-injection
corollary, and a proof\cite{tao}. The dimension lower bound is its
application to a locally full-rank empirical moment map.
\item \textbf{Set-summary dimensional obstructions are prior work.}
Wagstaff et al. make continuity versus discontinuous single-real
encodings explicit and prove continuous latent-dimension obstructions
for sum representations of set functions\cite{wagstaff}. Their
Theorem 4.1 proof uses an injection from an open ordered-set domain
into the latent space. Our specialization concerns a fixed finite
quadratic query family, and permits any continuous summary, including
nonadditive and order-dependent ones. The dimension obstruction and
the need to exclude pathological encodings are not new principles.
\end{itemize}

The defensible project contribution is the precise specialization:
the $q+2p$ count after polynomial dependencies and labels are charged;
the $q+p$ additional cost given $S_y$; a proof for sufficiently large
fixed equal-weight samples rather than only arbitrary weighted
measures; and the finite-query accuracy metric that yields a bit
packing bound. This package is presented as a precise boundary for the project's
historical summary, with classical ingredients credited and no claim of
priority. An algorithmic or learning-theoretic claim needs a separate result.
